Aus der Zustandsgleichung $pV = nRT$ folgt
\begin{align}
 V &= \frac{nR}{p}\cdot T
\end{align}
\begin{abcliste}
\abc Linie 1 ist waagrecht: $V$ bleibt gleich, also eine $\result{\text{Isochore}}$. Linie 3 ist senkrecht: $T$ bleibt gleich, also eine $\result{\text{Isotherme}}$. Die Linien 2 und 4 sind Geraden durch den Ursprung, $V \propto T$: Nach der Gleichung oben ist dann $p$ konstant, beide sind $\result{\text{Isobaren}}$.

\abc Die Steigung einer Isobaren ist $\frac{nR}{p}$: Je grösser der Druck, desto flacher die Gerade. In den Einheiten des Diagramms ist mit $T_0 = \frac{p_0V_0}{nR}$
\begin{align}
 \frac{V}{V_0} &= \frac{p_0}{p}\cdot\frac{T}{T_0}
\end{align}
die Steigung also $p_0/p$. Linie 2 steigt um $2$, Linie 4 um $\frac12$:
\begin{align}
 \text{Linie 2: } p &= \frac{p_0}{2} & \text{Linie 4: } p &= 2\,p_0
\end{align}
Auf der flacheren $\result{\text{Linie 4}}$ ist der Druck grösser, viermal so gross wie auf Linie 2. Anschaulich: Bei derselben Temperatur, etwa $T = 4\,T_0$, hat das Gas auf Linie 4 nur das Volumen $2\,V_0$, auf Linie 2 aber $8\,V_0$. Dasselbe Gas bei derselben Temperatur in einem kleineren Volumen hat den grösseren Druck.

\abc Auf Linie 1 ist $\result{V = 5\,V_0}$, auf Linie 3 ist $\result{T = 6\,T_0}$.
\end{abcliste}
